% ============================================================
% Paper 55 (DE): Faserbündel und charakteristische Klassen
% Autor: Michael Fuhrmann
% Projekt: Specialist Mathematics, Build 139
% Zuletzt geändert: 2026-03-12
% ============================================================
\documentclass[12pt,a4paper]{amsart}
\usepackage{amsmath,amssymb,amsthm}
\usepackage{mathtools}
\usepackage[utf8]{inputenc}
\usepackage[T1]{fontenc}
\usepackage{hyperref}
\usepackage{enumitem,booktabs,array}
\usepackage{geometry}
\usepackage{slashed}
\geometry{margin=2.5cm}

%% Theorem-Umgebungen (deutsch)
\newtheorem{theorem}{Satz}[section]
\newtheorem{lemma}[theorem]{Lemma}
\newtheorem{corollary}[theorem]{Korollar}
\newtheorem{proposition}[theorem]{Proposition}
\newtheorem{conjecture}[theorem]{Vermutung}
\theoremstyle{definition}
\newtheorem{definition}[theorem]{Definition}
\newtheorem{example}[theorem]{Beispiel}
\newtheorem{remark}[theorem]{Bemerkung}

\author{Michael Fuhrmann}
\address{Specialist Mathematics Project, Build 139}
\email{specialist-maths@localhost}
\date{\today}

\begin{document}
\title{Faserbündel und charakteristische Klassen:\\
Von der Topologie zur Eichtheorie}

\begin{abstract}
Faserbündel bilden den korrekten geometrischen Rahmen für die moderne Differenzialgeometrie,
Eichtheorie und mathematische Physik. Dieses Paper entwickelt die Theorie systematisch,
beginnend mit der grundlegenden Definition von Faserbündeln, Vektorbündeln und
Hauptfaserbündeln mit ihren Übergangsfunktionen und Strukturgruppen. Anschließend
werden Zusammenhänge und Krümmungsformen behandelt, die den Begriff des Paralleltransports
kodieren und über den Chern--Weil-Homomorphismus (bewiesen) zu charakteristischen Klassen
führen. Die zentralen Objekte sind die Chern-Klassen komplexer Vektorbündel, die
Pontryagin-Klassen reeller Bündel und die Euler-Klasse. Der berühmte
Atiyah--Singer-Indexsatz (bewiesen, 1963) vereint diese Ideen, indem er den analytischen
Index eines elliptischen Differenzialoperators als topologisches Integral charakteristischer
Klassen ausdrückt. Weiterhin behandeln wir $K$-Theorie und Bott-Periodizität (bewiesen),
die den natürlichen algebraischen Rahmen für charakteristische Klassen liefern, und
schließen mit physikalischen Anwendungen in Yang--Mills-Theorie, Instantonen und der
Berry-Phase in der Quantenmechanik. Durchgehend wird sorgfältig zwischen bewiesenen
Sätzen und offenen Vermutungen unterschieden.
\end{abstract}

\maketitle

\tableofcontents

% ============================================================
\section{Einleitung}
% ============================================================

Das Konzept eines Faserbündels stellt eine der tiefgreifendsten Verallgemeinerungen
kartesischer Produkträume in der modernen Mathematik dar. Während ein kartesisches
Produkt $B \times F$ die einfachste Weise bietet, zwei Räume zu kombinieren, erfordern
viele natürliche geometrische und physikalische Konstruktionen Räume, die lokal wie
Produkte aussehen, global jedoch ``verdreht'' sind. Das Möbiusband --- ein
Papierstreifen, der vor dem Verkleben der Enden um eine halbe Drehung verdreht wird
--- verdeutlicht dieses Phänomen exemplarisch: Es ist ein Linienbündel über dem Kreis
$S^1$, das global nicht trivializiert werden kann.

Faserbündel wurden von Whitney~\cite{whitney1935}, Seifert und Hopf in den 1930er
Jahren eingeführt; die grundlegende Theorie wurde von Steenrod~\cite{steenrod1951}
entwickelt. Ihre Bedeutung in der Mathematik erstreckt sich über algebraische
Topologie, Differenzialgeometrie und algebraische Geometrie. In der Physik sind
Faserbündel unverzichtbar: Das elektromagnetische Feld ist ein Zusammenhang auf
einem $U(1)$-Hauptfaserbündel über der Minkowski-Raumzeit, und das
Standardmodell der Teilchenphysik ist vollständig in der Sprache der
$SU(3) \times SU(2) \times U(1)$-Eichtheorien formuliert --- allesamt
Yang--Mills-Theorien auf Hauptfaserbündeln.

\subsection*{Überblick}
Abschnitt~\ref{sec:buendel} definiert Faserbündel in vollem Umfang, einschließlich
lokaler Trivializierungen und Übergangsfunktionen. Abschnitt~\ref{sec:vektor}
spezialisiert auf Vektorbündel und ihre Schnitte. Abschnitt~\ref{sec:haupt}
führt Hauptfaserbündel und Eichtransformationen ein. Zusammenhänge und Krümmung
sind Gegenstand von Abschnitt~\ref{sec:zusammenhaenge}. Charakteristische
Klassen --- Chern, Pontryagin und Euler --- werden in Abschnitt~\ref{sec:charakteristisch}
konstruiert. Der Chern--Weil-Homomorphismus wird in Abschnitt~\ref{sec:chernweil}
bewiesen. Der Atiyah--Singer-Indexsatz ist Inhalt von Abschnitt~\ref{sec:indexsatz}.
Physikalische Anwendungen behandelt Abschnitt~\ref{sec:physik}, und $K$-Theorie
beschließt den Text in Abschnitt~\ref{sec:ktheorie}.

\subsection*{Notation}
Im gesamten Text bezeichnet $M$ eine glatte $n$-dimensionale Mannigfaltigkeit
(kompakt, sofern nicht anders angegeben), $G$ eine Lie-Gruppe mit Lie-Algebra
$\mathfrak{g}$, und $H^*(X;\mathbb{Z})$ die singuläre Kohomologie eines
topologischen Raumes $X$ mit ganzzahligen Koeffizienten. Das Symbol
$\Omega^k(M, \mathfrak{g})$ bezeichnet den Raum der $\mathfrak{g}$-wertigen
Differenzial-$k$-Formen auf $M$.

% ============================================================
\section{Faserbündel}
\label{sec:buendel}
% ============================================================

\begin{definition}[Faserbündel]
Ein \emph{Faserbündel} besteht aus einem Tupel $(E, B, F, \pi, G)$, wobei
\begin{itemize}[noitemsep]
  \item $E$ der \emph{Gesamtraum} ist,
  \item $B$ der \emph{Basisraum} ist,
  \item $F$ die \emph{typische Faser} ist,
  \item $\pi \colon E \to B$ eine stetige surjektive Abbildung, die
        \emph{Projektion}, ist,
  \item $G$ die \emph{Strukturgruppe} ist, eine topologische Gruppe, die
        stetig von links auf $F$ wirkt,
\end{itemize}
sodass folgende \emph{lokale Trivialitätsbedingung} erfüllt ist: Es existiert
eine offene Überdeckung $\{U_\alpha\}_{\alpha \in \mathcal{A}}$ von $B$ und
Homöomorphismen
\[
  \varphi_\alpha \colon \pi^{-1}(U_\alpha) \xrightarrow{\;\sim\;} U_\alpha \times F,
\]
mit $\mathrm{pr}_1 \circ \varphi_\alpha = \pi$ (das Diagramm kommutiert).
\end{definition}

Die Struktur eines Faserbündels wird durch die \emph{Übergangsfunktionen} kodiert.
Auf Überschneidungen $U_\alpha \cap U_\beta \neq \emptyset$ definieren wir:
\[
  g_{\alpha\beta} \colon U_\alpha \cap U_\beta \to G, \quad
  g_{\alpha\beta}(b) = \varphi_\alpha|_b \circ (\varphi_\beta|_b)^{-1}.
\]
Diese Übergangsfunktionen erfüllen die \emph{Kozyklus-Bedingungen}:
\begin{align*}
  g_{\alpha\alpha}(b) &= e \quad \text{(Einselement in } G\text{)},\\
  g_{\alpha\beta}(b) \cdot g_{\beta\alpha}(b) &= e,\\
  g_{\alpha\beta}(b) \cdot g_{\beta\gamma}(b) \cdot g_{\gamma\alpha}(b) &= e
  \quad \text{für alle } b \in U_\alpha \cap U_\beta \cap U_\gamma.
\end{align*}

\begin{theorem}[Rekonstruktionssatz]
\label{thm:rekonstruktion}
Gegeben seien ein Basisraum $B$, eine Faser $F$, eine Strukturgruppe $G$, die
auf $F$ wirkt, eine offene Überdeckung $\{U_\alpha\}$ von $B$ und stetige
Abbildungen $g_{\alpha\beta} \colon U_\alpha \cap U_\beta \to G$, die die
Kozyklus-Bedingungen erfüllen. Dann existiert (bis auf Isomorphie von
Faserbündeln) ein eindeutiges Faserbündel $(E, B, F, \pi, G)$ mit diesen
Übergangsfunktionen.
\end{theorem}

\begin{proof}
Konstruiere $E$ als Quotient
$\bigsqcup_\alpha (U_\alpha \times F) / {\sim}$,
wobei $(b, f) \in U_\alpha \times F$ mit $(b, g_{\alpha\beta}(b) \cdot f)
\in U_\beta \times F$ identifiziert wird. Die Kozyklus-Bedingungen stellen
sicher, dass dies eine Äquivalenzrelation ist; $\pi$ wird durch Projektion auf
den ersten Faktor induziert.
\end{proof}

\begin{example}[Möbiusband]
Das Möbiusband ist ein reelles Linienbündel vom Rang $1$ über $S^1$ mit
Strukturgruppe $\mathbb{Z}_2 = \{+1, -1\}$. Wir überdecken $S^1$ durch zwei
offene Bögen $U_0$ und $U_1$; die Übergangsfunktion ist $g_{01}(b) = -1$ für
$b$ in einer Zusammenhangskomponente von $U_0 \cap U_1$ und $+1$ für die andere.
Dies ergibt ein nicht-orientierbares (daher nicht-triviales) Bündel: Es gibt
keine global konsistente Auszeichnung einer ``oberen'' Seite des Bandes.
\end{example}

\begin{example}[Hopf-Faserung]
Die \emph{Hopf-Faserung} $\eta \colon S^3 \to S^2$ ist definiert durch
\[
  \eta(z_1, z_2) = [z_1 : z_2] \in \mathbb{CP}^1 \cong S^2,
\]
wobei $(z_1, z_2) \in S^3 \subset \mathbb{C}^2$. Die Faser ist $S^1 \cong U(1)$,
und die Strukturgruppe ist $U(1)$. Die Hopf-Faserung ist das prototypische
nicht-triviale $U(1)$-Hauptfaserbündel über $S^2$ mit erster Chern-Zahl $c_1 = 1$.
\end{example}

\begin{remark}
Ein Faserbündel ist genau dann \emph{trivial}, wenn es zum Produktbündel
$B \times F \to B$ isomorph ist. Äquivalent dazu können alle Übergangsfunktionen
als konstante Abbildungen auf das Einselement $e \in G$ gewählt werden.
\end{remark}

% ============================================================
\section{Vektorbündel}
\label{sec:vektor}
% ============================================================

\begin{definition}[Vektorbündel]
Ein \emph{reelles (bzw.\ komplexes) Vektorbündel vom Rang $k$} über $B$ ist ein
Faserbündel $(E, B, \mathbb{R}^k, \pi, GL(k,\mathbb{R}))$ (bzw.\ mit Faser
$\mathbb{C}^k$ und Strukturgruppe $GL(k,\mathbb{C})$), sodass jede Faser
$\pi^{-1}(b)$ die Struktur eines $k$-dimensionalen reellen (bzw.\ komplexen)
Vektorraums trägt und die lokalen Trivializierungen faserweise lineare
Isomorphismen sind.
\end{definition}

\begin{definition}[Tangentialbündel]
Sei $M$ eine glatte $n$-Mannigfaltigkeit. Das \emph{Tangentialbündel} $TM$ ist
die disjunkte Vereinigung $TM = \bigsqcup_{p \in M} T_pM$ aller Tangentialräume,
ausgestattet mit der natürlichen Projektion $\pi \colon TM \to M$,
$v \mapsto p$ für $v \in T_pM$. Dies ist ein reelles Vektorbündel vom Rang $n$
über $M$.
\end{definition}

\begin{example}[Kotangential- und Normalbündel]
Das \emph{Kotangentialbündel} $T^*M$ ist das duale Bündel $(TM)^*$; seine
Schnitte sind Differenzial-$1$-Formen. Für eine Einbettung $M \hookrightarrow N$
glatter Mannigfaltigkeiten besteht das \emph{Normalbündel} $\nu$ aus den
(bezüglich einer gewählten Riemannschen Metrik) zu $TM$ senkrechten Vektoren
in $TN$, was die kurze exakte Sequenz
$0 \to TM \to TN|_M \to \nu \to 0$
liefert.
\end{example}

\begin{definition}[Schnitt (Section)]
Ein \emph{Schnitt} eines Vektorbündels $\pi \colon E \to B$ ist eine stetige
Abbildung $s \colon B \to E$ mit $\pi \circ s = \mathrm{id}_B$. Der Raum der
glatten Schnitte wird mit $\Gamma(E)$ bezeichnet. Für das Tangentialbündel sind
Schnitte genau die \emph{Vektorfelder} auf $M$.
\end{definition}

\begin{theorem}[Igelbürstenlemma, bewiesen]
\label{thm:igel}
Das Tangentialbündel $TS^n$ besitzt genau dann einen nirgendwo verschwindenden
Schnitt, wenn $n$ ungerade ist.
\end{theorem}

Insbesondere muss jedes Vektorfeld auf $S^2$ irgendwo verschwinden --- man
kann einen Igel nicht ohne Kuhle kämmen.

\begin{definition}[Globale Trivializierung und stabil triviale Bündel]
Ein Vektorbündel $E$ ist \emph{trivializierbar} (global trivial), wenn es
$k$ Schnitte $s_1, \ldots, s_k \in \Gamma(E)$ gibt, die in jeder Faser eine
Basis bilden. Ein Bündel heißt \emph{stabil trivial}, wenn
$E \oplus \varepsilon^m \cong \varepsilon^{k+m}$ für ein triviales Bündel
$\varepsilon^m$ gilt.
\end{definition}

Operationen auf Vektorbündeln umfassen direkte Summe $E \oplus F$, Tensorprodukt
$E \otimes F$, das duale Bündel $E^*$ und äußere Potenzen $\Lambda^k E$. Das
\emph{Spaltungsprinzip} (Splitting Principle) besagt: Jede Formel für
charakteristische Klassen, die für direkte Summen von Linienbündeln gilt,
überträgt sich auf allgemeine Bündel.

% ============================================================
\section{Hauptfaserbündel}
\label{sec:haupt}
% ============================================================

Hauptfaserbündel sind die grundlegenden Objekte der Eichtheorie. Ein Vektorbündel
lässt sich stets aus einem zugehörigen Hauptfaserbündel zurückgewinnen, was
Hauptfaserbündel zu den ``Masterobjekten'' macht.

\begin{definition}[Haupt-$G$-Bündel]
Ein \emph{Haupt-$G$-Bündel} ist ein Faserbündel $\pi \colon P \to B$ zusammen
mit einer stetigen freien Rechts-$G$-Wirkung $P \times G \to P$,
$(p, g) \mapsto p \cdot g$, sodass:
\begin{enumerate}[noitemsep]
  \item $\pi(p \cdot g) = \pi(p)$ für alle $p \in P$, $g \in G$
        (die Wirkung erhält Fasern),
  \item die Wirkung auf jeder Faser frei und transitiv ist
        (jede Faser ist als $G$-Torseur isomorph zu $G$),
  \item lokale Trivializierungen $\varphi_\alpha \colon \pi^{-1}(U_\alpha) \to
        U_\alpha \times G$ $G$-äquivariant sind.
\end{enumerate}
\end{definition}

\begin{example}
Das \emph{Rahmenbündel} $\mathrm{Fr}(M)$ einer glatten Mannigfaltigkeit $M$ ist
das $GL(n,\mathbb{R})$-Hauptfaserbündel, dessen Faser über $p \in M$ aus allen
geordneten Basen (Rahmen) von $T_pM$ besteht. Durch Reduktion auf orthonormale
Rahmen erhält man ein $O(n)$-Bündel; bei Wahl einer Orientierung ein $SO(n)$-Bündel.
\end{example}

\begin{definition}[Assoziiertes Vektorbündel]
Gegeben ein $G$-Hauptfaserbündel $P \to B$ und eine Darstellung
$\rho \colon G \to GL(V)$ auf einem Vektorraum $V$, ist das
\emph{assoziierte Vektorbündel}
\[
  E = P \times_G V := (P \times V) / G,
\]
wobei die Rechts-$G$-Wirkung durch $(p, v) \cdot g = (p \cdot g, \rho(g^{-1}) v)$
gegeben ist.
\end{definition}

\begin{definition}[Eichtransformation]
Eine \emph{Eichtransformation} eines $G$-Hauptfaserbündels $P \to B$ ist ein
Bündelautomorphismus $\phi \colon P \to P$, der die Identität auf $B$ überdeckt
und $G$-äquivariant ist: $\phi(p \cdot g) = \phi(p) \cdot g$. Eichtransformationen
bilden die \emph{Eichgruppe} $\mathcal{G} = \Gamma(\mathrm{Ad}(P))$.
\end{definition}

\begin{proposition}
Jedes $G$-Hauptfaserbündel $P \to B$ ist bis auf Isomorphismus durch seine
Übergangsfunktionen $g_{\alpha\beta} \colon U_\alpha \cap U_\beta \to G$
bestimmt, die eine \v{C}ech-Kohomologieklasse in $\check{H}^1(B, G)$
repräsentieren.
\end{proposition}

% ============================================================
\section{Zusammenhänge auf Bündeln}
\label{sec:zusammenhaenge}
% ============================================================

Ein \emph{Zusammenhang} (Konnexion) auf einem Faserbündel liefert eine
systematische Methode, Schnitte zu differenzieren und Fasern entlang von Kurven
im Basisraum parallel zu transportieren.

\begin{definition}[Zusammenhang auf einem Hauptfaserbündel]
Ein \emph{Zusammenhang} auf einem $G$-Hauptfaserbündel $P \to B$ ist eine
$\mathfrak{g}$-wertige $1$-Form $A \in \Omega^1(P, \mathfrak{g})$
mit den Eigenschaften:
\begin{enumerate}[noitemsep]
  \item (Äquivarianz)\; $R_g^* A = \mathrm{Ad}(g^{-1}) A$ für alle $g \in G$,
  \item (Reproduktion)\; $A(\xi^*) = \xi$ für alle $\xi \in \mathfrak{g}$,
        wobei $\xi^*$ das zu $\xi$ gehörende fundamentale Vektorfeld ist.
\end{enumerate}
\end{definition}

Ein Zusammenhang $A$ ist äquivalent durch eine Wahl horizontaler Unterräume
$H_p \subset T_p P$ gegeben, die glatt variieren und
$T_pP = H_p \oplus V_p$ erfüllen (mit $V_p = \ker(d\pi)_p$ dem vertikalen
Unterraum), sowie $G$-äquivariant sind.

\begin{definition}[Krümmungsform]
Die \emph{Krümmung} eines Zusammenhangs $A$ ist die $\mathfrak{g}$-wertige
$2$-Form
\[
  F_A = dA + \tfrac{1}{2}[A \wedge A] = dA + A \wedge A \in \Omega^2(P, \mathfrak{g}).
\]
In lokalen Koordinaten (Eichpotenzial $A_\mu \, dx^\mu$ auf $B$):
\[
  F_{\mu\nu} = \partial_\mu A_\nu - \partial_\nu A_\mu + [A_\mu, A_\nu].
\]
Für abelsche Gruppen $G$ (z.B.\ $U(1)$) verschwindet die Klammer: $F = dA$.
\end{definition}

\begin{theorem}[Bianchi-Identität, bewiesen]
\label{thm:bianchi}
Für jeden Zusammenhang $A$ auf einem Hauptfaserbündel gilt $D_A F_A = 0$,
wobei $D_A = d + [A, \cdot]$ die kovariante äußere Ableitung ist. In lokalen
Koordinaten:
$\partial_\lambda F_{\mu\nu} + \partial_\mu F_{\nu\lambda} + \partial_\nu F_{\lambda\mu}
+ [A_\lambda, F_{\mu\nu}] + \cdots = 0.$
\end{theorem}

\begin{definition}[Paralleltransport]
Gegeben eine glatte Kurve $\gamma \colon [0,1] \to B$ und ein Zusammenhang $A$,
ist der \emph{Paralleltransport} $P_\gamma \colon \pi^{-1}(\gamma(0)) \to
\pi^{-1}(\gamma(1))$ durch das Anheben von $\gamma$ zu einer horizontalen
Kurve in $P$ definiert --- einer Kurve, deren Tangentialvektoren in der
horizontalen Distribution $H$ liegen. Bei Vektorbündeln ergibt dies einen
linearen Isomorphismus zwischen den Fasern.
\end{definition}

\begin{definition}[Holonomie]
Für eine geschlossene Schleife $\gamma$ mit Basispunkt $b_0 \in B$ liefert
der Paralleltransport ein Element $\mathrm{Hol}(A, \gamma) \in G$, die
\emph{Holonomie}. Für abelsche Gruppen gilt:
$\mathrm{Hol}(A, \gamma) = \exp\!\Bigl(\oint_\gamma A\Bigr)
= \exp\!\Bigl(\iint_\Sigma F_A\Bigr),$
wobei $\Sigma$ eine beliebige Fläche mit $\partial\Sigma = \gamma$ ist
(Satz von Stokes).
\end{definition}

\begin{example}[Berry-Phase]
In der Quantenmechanik betrachten wir einen Hamiltonoperator $H(\lambda)$, der
von Parametern $\lambda \in B$ abhängt, und einen Eigenzustand $\psi(\lambda)$,
der sich adiabatisch entwickelt. Wenn die Parameter eine geschlossene Schleife
$\gamma \subset B$ beschreiben, erwirbt der Zustand eine Phase $e^{i\varphi_B}$
mit
$\varphi_B = \oint_\gamma \langle \psi | d\psi \rangle,$
die \emph{Berry-Phase}~\cite{berry1984}. Dies ist genau die Holonomie des
natürlichen $U(1)$-Zusammenhangs auf dem Linienbündel über $B$, das durch die
Eigenzustände gebildet wird.
\end{example}

% ============================================================
\section{Charakteristische Klassen}
\label{sec:charakteristisch}
% ============================================================

Charakteristische Klassen sind Kohomologieklassen, die natürlich zu Vektorbündeln
gehören. Sie liefern (in vielen Fällen vollständige) topologische Hindernisse
gegen Trivializierbarkeit, Orientierbarkeit und andere geometrische Eigenschaften.

\subsection{Chern-Klassen}

\begin{definition}[Chern-Klassen]
Sei $E \to B$ ein komplexes Vektorbündel vom Rang $k$. Die \emph{Chern-Klassen}
$c_i(E) \in H^{2i}(B; \mathbb{Z})$, $0 \le i \le k$, sind die eindeutig
bestimmten Kohomologieklassen, die folgende Axiome erfüllen:
\begin{enumerate}[noitemsep]
  \item \textbf{Natürlichkeit}: $f^* c_i(E) = c_i(f^*E)$ für stetige
        Abbildungen $f \colon B' \to B$.
  \item \textbf{Whitney-Summenformel}: $c(E \oplus F) = c(E) \cup c(F)$,
        wobei $c(E) = 1 + c_1(E) + \cdots + c_k(E)$ die totale Chern-Klasse ist.
  \item \textbf{Normierung}: Für das tautologische Linienbündel $\gamma^1$ über
        $\mathbb{CP}^1$ gilt $c_1(\gamma^1) = -[\omega]$, das Negative des
        kanonischen Erzeugers von $H^2(\mathbb{CP}^1;\mathbb{Z}) \cong \mathbb{Z}$.
  \item $c_0(E) = 1$ und $c_i(E) = 0$ für $i > \mathrm{rang}(E)$.
\end{enumerate}
\end{definition}

\begin{example}
Die Hopf-Faserung $\eta \colon S^3 \to S^2$ entspricht dem tautologischen
Linienbündel $\mathcal{O}(-1)$ über $\mathbb{CP}^1 \cong S^2$ mit $c_1 = -1$.
Das duale Bündel $\mathcal{O}(1)$ hat $c_1 = +1$.
\end{example}

\subsection{Pontryagin-Klassen}

\begin{definition}[Pontryagin-Klassen]
Für ein reelles Vektorbündel $E \to B$ vom Rang $k$ sind die \emph{Pontryagin-Klassen}
$p_i(E) \in H^{4i}(B; \mathbb{Z})$ durch Komplexifizierung definiert:
\[
  p_i(E) = (-1)^i c_{2i}(E \otimes_\mathbb{R} \mathbb{C}).
\]
Die totale Pontryagin-Klasse ist $p(E) = 1 + p_1(E) + p_2(E) + \cdots$.
\end{definition}

\begin{remark}
Pontryagin-Klassen liegen in $4i$-fachen Kohomologiegruppen (nicht $2i$-fachen
wie Chern-Klassen), da die Komplexifizierung eines reellen Bündels vom Rang $k$
ein komplexes Bündel vom Rang $k$ ergibt und die ungeraden Chern-Klassen der
Komplexifizierung verschwinden.
\end{remark}

\begin{theorem}[Pontryagin-Signatursatz, bewiesen]
\label{thm:signatur}
Für eine kompakte orientierte $4k$-Mannigfaltigkeit $M$ erfüllt die Signatur
$\sigma(M) = \mathrm{Vorzeichen}(H^{2k}(M;\mathbb{R}))$:
$\sigma(M) = \langle L(M), [M] \rangle$,
wobei $L(M)$ das Hirzebruch-$L$-Polynom in den Pontryagin-Klassen ist.
\end{theorem}

\subsection{Euler-Klasse}

\begin{definition}[Euler-Klasse]
Für ein orientiertes reelles Vektorbündel $E \to B$ vom Rang $k$ ist die
\emph{Euler-Klasse} $e(E) \in H^k(B;\mathbb{Z})$ definiert als
Selbstschnitt-Klasse des Nullschnitts von $E$, oder äquivalent als das
Hindernis für einen nirgendwo verschwindenden Schnitt.
\end{definition}

\begin{theorem}[Euler-Klasse und Euler-Charakteristik, bewiesen]
Für eine kompakte orientierte $n$-Mannigfaltigkeit $M$ gilt:
$\chi(M) = \langle e(TM), [M] \rangle$,
der Grad der Euler-Klasse des Tangentialbündels.
\end{theorem}

% ============================================================
\section{Chern--Weil-Theorie}
\label{sec:chernweil}
% ============================================================

Der Chern--Weil-Homomorphismus ist die Schlüsselkonstruktion, die aus der Krümmung
eines Zusammenhangs de Rham-Kohomologieklassen (und damit charakteristische Klassen)
erzeugt. Er bildet eine elegante Brücke zwischen Differenzialgeometrie und
algebraischer Topologie.

\begin{definition}[Invariantes Polynom]
Ein Polynom $P \colon \mathfrak{g} \to \mathbb{R}$ heißt
\emph{$\mathrm{Ad}$-invariant}, wenn
$P(\mathrm{Ad}(g) X) = P(X)$ für alle $g \in G$ und $X \in \mathfrak{g}$ gilt.
Der Ring aller solchen Polynome wird mit $I^*(G) = \mathbb{R}[\mathfrak{g}]^G$
bezeichnet.
\end{definition}

\begin{theorem}[Chern--Weil-Homomorphismus, bewiesen]
\label{thm:chernweil}
Sei $P \to B$ ein $G$-Hauptfaserbündel mit Zusammenhang $A$ und Krümmungsform
$F_A \in \Omega^2(P, \mathfrak{g})$. Für jedes $\mathrm{Ad}$-invariante Polynom
$P \in I^k(G)$ (homogen vom Grad $k$) ist die Differenzialform
\[
  P(F_A) \in \Omega^{2k}(B)
\]
geschlossen: $d P(F_A) = 0$, und ihre de Rham-Kohomologieklasse $[P(F_A)] \in
H^{2k}_{\mathrm{dR}}(B)$ ist von der Wahl des Zusammenhangs $A$ unabhängig.
Der resultierende Ringhomomorphismus
\[
  \mathrm{cw} \colon I^*(G) \longrightarrow H^*_{\mathrm{dR}}(B)
\]
ist der \textup{Chern--Weil-Homomorphismus}.
\end{theorem}

\begin{proof}[Beweisidee]
\textbf{Geschlossenheit}: Mit der Bianchi-Identität $D_A F_A = 0$ und der
$\mathrm{Ad}$-Invarianz von $P$ zeigt eine direkte Rechnung $d P(F_A) = 0$.

\textbf{Unabhängigkeit vom Zusammenhang}: Seien $A_0$ und $A_1$ zwei
Zusammenhänge. Betrachte die Familie $A_t = (1-t)A_0 + tA_1$ auf $P \times [0,1]$.
Man berechnet:
\[
  P(F_{A_1}) - P(F_{A_0}) = d\!\int_0^1 Q(A_1 - A_0, F_{A_t})\, dt,
\]
wobei $Q$ die aus $P$ abgeleitete Transgressionsform ist. Daher gilt
$[P(F_{A_1})] = [P(F_{A_0})]$ in $H^*_{\mathrm{dR}}(B)$.
\end{proof}

\begin{corollary}
Die totale Chern-Klasse eines komplexen Vektorbündels $E$ mit Krümmung $F_A$ ist:
\[
  c(E) = \left[\det\!\left(I + \tfrac{i}{2\pi} F_A\right)\right]
  \in H^{\mathrm{gerade}}_{\mathrm{dR}}(B),
\]
die einzelnen Chern-Klassen $c_k(E)$ sind die Grad-$2k$-Anteile dieser Determinantenentwicklung.
Hinweis: Die Spur-Formel $\mathrm{tr}((i/2\pi)\,F_A)^k/k!$ liefert die \emph{Chern-Charakter}-Komponente
$\mathrm{ch}_k(E)$, nicht die Chern-Klasse $c_k(E)$.
\end{corollary}

Die Stärke des Chern--Weil-Satzes liegt darin, dass er gleichzeitig die Existenz
charakteristischer Klassen nachweist und eine explizite Krümmungsformel zur
Berechnung liefert.

% ============================================================
\section{Der Atiyah--Singer-Indexsatz}
\label{sec:indexsatz}
% ============================================================

Der Atiyah--Singer-Indexsatz ist eines der tiefgründigsten Ergebnisse der
Mathematik des 20.\ Jahrhunderts. Er vereint Analysis (Spektraltheorie
von Differenzialoperatoren), Topologie (charakteristische Klassen) und
Geometrie (Riemannsche Mannigfaltigkeiten).

\begin{definition}[Elliptischer Operator und analytischer Index]
Seien $E, F \to M$ Vektorbündel über einer kompakten Mannigfaltigkeit $M$ und
$D \colon \Gamma(E) \to \Gamma(F)$ ein linearer Differenzialoperator. $D$
heißt \emph{elliptisch}, wenn sein Hauptsymbol
$\sigma(D)(x, \xi) \colon E_x \to F_x$ für alle $x \in M$ und $\xi \neq 0$
ein Isomorphismus ist. Der \emph{analytische Index} eines elliptischen
Operators ist
\[
  \mathrm{ind}(D) = \dim \ker D - \dim \ker D^*.
\]
Durch Elliptizität und Kompaktheit sind beide Dimensionen endlich.
\end{definition}

\begin{theorem}[Atiyah--Singer-Indexsatz, bewiesen 1963]
\label{thm:indexsatz}
Sei $D \colon \Gamma(E) \to \Gamma(F)$ ein elliptischer Differenzialoperator
auf einer kompakten orientierbaren $n$-Mannigfaltigkeit $M$ ohne Rand. Dann gilt:
\[
  \mathrm{ind}(D) = \int_M \mathrm{ch}(\sigma(D)) \cdot \mathrm{Td}(TM \otimes \mathbb{C}),
\]
wobei $\mathrm{ch}(\sigma(D))$ der Chern-Charakter der Symbolklasse in $K(T^*M)$
und $\mathrm{Td}(TM \otimes \mathbb{C})$ die Todd-Klasse des komplexifizierten
Tangentialbündels ist.
\end{theorem}

\begin{remark}
Der Satz wurde erstmals von Atiyah und Singer im Jahr 1963 bewiesen~\cite{atiyahsinger1963},
mit späteren Vereinfachungen durch $K$-Theorie und Wärme-Kern-Methoden. Atiyah
erhielt 1966 die Fields-Medaille, unter anderem für dieses Werk. Der Satz
verallgemeinert folgende klassische Ergebnisse:
\begin{itemize}[noitemsep]
  \item \textbf{Gauss--Bonnet}: $\chi(M) = \frac{1}{2\pi}\int_M K\, dA$
        (für den de Rham-Operator $d + d^*$).
  \item \textbf{Hirzebruch-Signatursatz}: $\sigma(M) = \langle L(p(M)), [M]\rangle$.
  \item \textbf{Riemann--Roch}: $\chi(\mathcal{O}(D)) = \deg D + 1 - g$
        (für den Dolbeault-Operator auf einer Riemannschen Fläche).
\end{itemize}
\end{remark}

\begin{corollary}[Dirac-Operator und $\hat{A}$-Geschlecht]
Für den Dirac-Operator $\slashed{D}$ auf einer kompakten Spin-Mannigfaltigkeit
$M$ der Dimension $4k$ gilt:
\[
  \mathrm{ind}(\slashed{D}) = \hat{A}(M) = \int_M \hat{A}(TM),
\]
wobei $\hat{A}$ das $\hat{A}$-Geschlecht (ein Polynom in den Pontryagin-Klassen)
ist. Insbesondere gilt $\hat{A}(M) \in \mathbb{Z}$.
\end{corollary}

% ============================================================
\section{Anwendungen in der Physik}
\label{sec:physik}
% ============================================================

\subsection{Yang--Mills-Theorie}

Der mathematische Rahmen von Eichtheorien ist genau die Theorie der Zusammenhänge
auf Hauptfaserbündeln. Das Yang--Mills-Wirkungsfunktional ist für Zusammenhänge
auf einem $G$-Hauptfaserbündel $P \to M$ über einer Riemannschen $4$-Mannigfaltigkeit
$(M, g)$ definiert:
\[
  S_{\mathrm{YM}}(A) = \frac{1}{2} \int_M \|F_A\|^2 \, d\mathrm{vol}_g
  = \frac{1}{2} \int_M \mathrm{tr}(F_A \wedge {*F_A}),
\]
wobei $*$ der Hodge-Stern-Operator ist. Die \emph{Yang--Mills-Gleichungen}
(Euler--Lagrange-Gleichungen von $S_{\mathrm{YM}}$) lauten:
\[
  D_A {*F_A} = 0 \quad \text{und} \quad D_A F_A = 0 \text{ (Bianchi)}.
\]

\begin{definition}[Instanton]
Ein Zusammenhang $A$ auf einem $G$-Hauptfaserbündel über $S^4$ heißt
\emph{Instanton} (oder \emph{selbstdualer Yang--Mills-Zusammenhang}), wenn
\[
  F_A = {*F_A} \quad (\text{selbstdual}),
\]
oder $F_A = -{*F_A}$ (antidual). Instantonen sind absolute Minima von
$S_{\mathrm{YM}}$ innerhalb eines festen topologischen Sektors, der durch die
\emph{Instanton-Zahl}
\[
  k = \frac{1}{8\pi^2} \int_{S^4} \mathrm{tr}(F_A \wedge F_A) \in \mathbb{Z}
\]
charakterisiert wird.
\end{definition}

\begin{theorem}[BPST-Instanton, bewiesen]
\label{thm:bpst}
Für jedes $k \in \mathbb{Z}$ existiert ein selbstdualer $SU(2)$-Zusammenhang
auf dem Hauptfaserbündel über $S^4$ mit Instanton-Zahl $k$. Für $k=1$ ist der
\emph{BPST-Instanton} (Belavin--Polyakov--Schwarz--Tyupkin, 1975) in singulärer
Eichung gegeben durch:
\[
  A_\mu^a = \frac{2}{g} \frac{\eta_{\mu\nu}^a \, x^\nu}{x^2 + \rho^2},
\]
wobei $\eta^a_{\mu\nu}$ die 't-Hooft-Symbole und $\rho > 0$ die Instanton-Größe sind.
\end{theorem}

\begin{remark}
Der Modulraum der $k$ Instantonen auf $S^4$ für $G = SU(2)$ hat Dimension
$8k - 3$. Für $k=1$ ergibt dies einen $5$-dimensionalen Modulraum (Position:
$4$ Parameter, Größe: $1$ Parameter). Instantonen haben tiefgreifende
Konsequenzen für nicht-perturbative Effekte in der Quantenfeldtheorie,
einschließlich des Tunnelns zwischen topologisch verschiedenen Vakua.
\end{remark}

\subsection{Quanten-Hall-Effekt und Chern-Zahlen}

Der ganzzahlige Quanten-Hall-Effekt liefert eine physikalische Realisierung
der ersten Chern-Zahl. Die Hall-Leitfähigkeit eines zweidimensionalen
Elektronensystems im Magnetfeld ist quantisiert als
\[
  \sigma_{xy} = \frac{e^2}{h} \cdot C_1,
\]
wobei $C_1 = \frac{1}{2\pi} \int_{\mathrm{BZ}} F \in \mathbb{Z}$ die erste
Chern-Zahl des Berry-Zusammenhangs über dem Brillouin-Zonen-Torus ist
(TKNN-Formel, Thouless--Kohmoto--Nightingale--den~Nijs, 1982).

% ============================================================
\section{$K$-Theorie}
\label{sec:ktheorie}
% ============================================================

Die $K$-Theorie liefert den natürlichen algebraischen Rahmen für die Untersuchung
von Vektorbündeln bis auf stabile Äquivalenz und verfeinert die Theorie der
charakteristischen Klassen.

\begin{definition}[$K$-Gruppe]
Die (komplexe, reduzierte) \emph{$K$-Gruppe} $K(X)$ eines kompakten
Hausdorff-Raumes $X$ ist die Grothendieck-Gruppe des Monoids der
Isomorphieklassen komplexer Vektorbündel über $X$ unter direkter Summe. Das
heißt, $K(X)$ besteht aus formalen Differenzen $[E] - [F]$ von Vektorbündeln
modulo der Relation $[E] - [F] = [E'] - [F']$ genau dann, wenn
$E \oplus F' \oplus \varepsilon^n \cong E' \oplus F \oplus \varepsilon^n$
für ein $n$ gilt.
\end{definition}

\begin{theorem}[Bott-Periodizität, bewiesen von Bott 1957]
\label{thm:bott}
Für jeden kompakten Hausdorff-Raum $X$ gibt es natürliche Isomorphismen
\[
  \tilde{K}(\Sigma^2 X) \cong \tilde{K}(X),
\]
wobei $\Sigma$ die reduzierte Einhängung bezeichnet. Äquivalent dazu ist das
$K$-Theorie-Spektrum $2$-periodisch: $K^{n+2}(X) \cong K^n(X)$.
Insbesondere gilt:
\[
  K(\mathrm{pt}) \cong \mathbb{Z}, \quad K^{-1}(\mathrm{pt}) \cong 0,
  \quad K^{-2}(\mathrm{pt}) \cong \mathbb{Z}, \quad \ldots
\]
\end{theorem}

\begin{remark}
Die Bott-Periodizität hat einen geometrischen Beweis: $\pi_{2k}(U) \cong \mathbb{Z}$
und $\pi_{2k+1}(U) \cong 0$ für alle $k \ge 0$ und die stabile unitäre Gruppe
$U = \varinjlim U(n)$. Die reelle Variante mit Periode $8$ (reelle Bott-Periodizität)
hängt mit den acht reellen Divisionsalgebren und der Periodizität der
Clifford-Algebren zusammen.
\end{remark}

\begin{theorem}[Atiyah--Hirzebruch-Spektralfolge, bewiesen]
\label{thm:ahss}
Für jeden endlichen CW-Komplex $X$ existiert eine Spektralfolge, die gegen
$K^*(X)$ konvergiert, mit $E_2$-Seite
$E_2^{p,q} \cong H^p(X; K^q(\mathrm{pt}))$.
\end{theorem}

\begin{corollary}
Der Chern-Charakter $\mathrm{ch} \colon K(X) \otimes \mathbb{Q} \xrightarrow{\sim}
H^{\mathrm{gerade}}(X; \mathbb{Q})$ ist ein Ringisomorphismus und liefert die
präzise Beziehung zwischen $K$-Theorie und rationaler Kohomologie.
\end{corollary}

\begin{theorem}[Atiyah--Singer via $K$-Theorie, bewiesen]
Der Atiyah--Singer-Indexsatz lässt sich als Aussage über $K$-Theorie
reformulieren: Die analytische Indexabbildung $\mathrm{ind}_a \colon K(T^*M)
\to \mathbb{Z}$ stimmt mit der topologischen Indexabbildung
$\mathrm{ind}_t \colon K(T^*M) \to \mathbb{Z}$ überein, die über den
Thom-Isomorphismus und das Vorwärts-Schieben in der $K$-Theorie definiert wird.
\end{theorem}

% ============================================================
\section{Literatur}
% ============================================================

\begin{thebibliography}{99}

\bibitem{milnorstasheff1974}
J.~W. Milnor und J.~D. Stasheff,
\textit{Characteristic Classes},
Annals of Mathematics Studies, Bd.~76,
Princeton University Press, Princeton, 1974.

\bibitem{husemoller1994}
D.~Husem\"{o}ller,
\textit{Fibre Bundles},
3.\ Aufl., Graduate Texts in Mathematics, Bd.~20,
Springer, New York, 1994.

\bibitem{atiyahsinger1963}
M.~F. Atiyah und I.~M. Singer,
,,The index of elliptic operators on compact manifolds,``
\textit{Bull.\ Amer.\ Math.\ Soc.}, \textbf{69} (1963), 422--433.

\bibitem{atiyahsinger1968}
M.~F. Atiyah und I.~M. Singer,
,,The index of elliptic operators: I,``
\textit{Ann.\ of Math.}, \textbf{87} (1968), 484--530.

\bibitem{botttu1982}
R.~Bott und L.~W. Tu,
\textit{Differential Forms in Algebraic Topology},
Graduate Texts in Mathematics, Bd.~82,
Springer, New York, 1982.

\bibitem{steenrod1951}
N.~Steenrod,
\textit{The Topology of Fibre Bundles},
Princeton Mathematical Series, Bd.~14,
Princeton University Press, Princeton, 1951.

\bibitem{berry1984}
M.~V. Berry,
,,Quantal phase factors accompanying adiabatic changes,``
\textit{Proc.\ R.\ Soc.\ Lond.\ A}, \textbf{392} (1984), 45--57.

\bibitem{whitney1935}
H.~Whitney,
,,Sphere spaces,``
\textit{Proc.\ Natl.\ Acad.\ Sci.\ USA}, \textbf{21} (1935), 462--468.

\bibitem{donaldson1983}
S.~K. Donaldson,
,,An application of gauge theory to the topology of $4$-manifolds,``
\textit{J.\ Differential Geom.}, \textbf{18} (1983), 269--316.

\bibitem{atiyah1967}
M.~F. Atiyah,
\textit{$K$-Theory},
Benjamin, New York, 1967.

\end{thebibliography}

\end{document}
